\documentclass[11pt,a4paper]{article}
\usepackage[utf8]{inputenc}
\usepackage[T1]{fontenc}
\usepackage[spanish]{babel}
\usepackage{amsmath,amssymb,mathtools}
\usepackage{geometry}
\usepackage{enumitem}
\usepackage{lmodern}
\usepackage{microtype}
\geometry{margin=2.4cm}
\setlength{\parindent}{0pt}
\setlength{\parskip}{0.45em}
\newcommand{\R}{\mathbb{R}}
\newcommand{\C}{\mathbb{C}}
\newcommand{\K}{\mathbb{K}}
\newcommand{\Ker}{\operatorname{Ker}}
\newcommand{\rg}{\operatorname{rg}}
\newcommand{\sg}{\operatorname{sg}}
\newcommand{\Span}{\operatorname{L}}
\begin{document}

\begin{center}
{\LARGE\bfseries Álgebra Lineal}\\[0.4em]
{\Large Tema 5: Formas Bilineales y Cuadráticas --- Soluciones}
\end{center}

\section{Ejercicio 1}
\subsection{Solución}

\[
f((x,y),(x',y'))
=
\begin{pmatrix}x&y\end{pmatrix}
\begin{pmatrix}
1&4\\
-3&0
\end{pmatrix}
\begin{pmatrix}x'\\y'\end{pmatrix}.
\]

Por tanto:

\[
\boxed{f\text{ es bilineal}.}
\]

\section{Ejercicio 2}
\subsection{Solución}

Una forma bilineal debe verificar:

\[
f((x,y),(0,0))=0.
\]

Pero:

\[
f((x,y),(0,0))=5x-3y.
\]

Por tanto:

\[
\boxed{f\text{ no es bilineal}.}
\]

\section{Ejercicio 3}
\subsection{Solución}

\[
f(0,0)=1\neq0.
\]

Por tanto:

\[
\boxed{f\text{ no es bilineal}.}
\]

\section{Ejercicio 4}
\subsection{Solución}

Para $\alpha,\beta\in\R$:

\[
f(\alpha p+\beta r,q)
=
\int_a^b(\alpha p+\beta r)q
=
\alpha f(p,q)+\beta f(r,q).
\]

Análogamente en la segunda componente.

Por tanto:

\[
\boxed{f\text{ es bilineal}.}
\]

\section{Ejercicio 5}
\subsection{Solución}

Definimos:

\[
f_s(u,v)
=
\frac12\left(f(u,v)+f(v,u)\right)
\]

y:

\[
f_a(u,v)
=
\frac12\left(f(u,v)-f(v,u)\right).
\]

Entonces:

\[
f_s(v,u)=f_s(u,v)
\]

y:

\[
f_a(v,u)=-f_a(u,v).
\]

Además:

\[
f_s(u,v)+f_a(u,v)=f(u,v).
\]

Por tanto:

\[
\boxed{f=f_s+f_a.}
\]

\section{Ejercicio 6}
\subsection{Solución}

La matriz es:

\[
\boxed{
A=
\begin{pmatrix}
1&0&-1\\
0&1&1\\
-1&1&2
\end{pmatrix}.
}
\]

\[
\rg A=2.
\]

Por tanto:

\[
\boxed{\rg f=2}
\]

y:

\[
AX=0
\]

da:

\[
x=z,\qquad y=-z.
\]

Así:

\[
\boxed{
N(f)=\Span\{(1,-1,1)\}.
}
\]

Del cambio de coordenadas:

\[
X'=CX,
\qquad
C=
\begin{pmatrix}
2&2&3\\
-1&0&-1\\
0&-1&-1
\end{pmatrix},
\]

se obtiene:

\[
X=C^{-1}X'.
\]

Por tanto:

\[
\boxed{
M_B(f)
=
(C^{-1})^tAC^{-1}.
}
\]

\section{Ejercicio 7}
\subsection{Solución}

Sea:

\[
M_B(f)=
\begin{pmatrix}
a_{11}&a_{12}\\
a_{12}&a_{22}
\end{pmatrix}.
\]

Como:

\[
f(x,x)=3,
\]

\[
a_{22}=3.
\]

Además:

\[
a_{22}+2a_{12}+a_{11}=8
\]

y:

\[
a_{22}+4a_{12}+4a_{11}=11.
\]

Resolviendo:

\[
a_{11}=-1,
\qquad
a_{12}=3.
\]

Por tanto:

\[
\boxed{
M_B(f)=
\begin{pmatrix}
-1&3\\
3&3
\end{pmatrix}.
}
\]

\section{Ejercicio 8}
\subsection{Solución}

Sean:

\[
p_1=2+x-x^2,
\quad
p_2=1+2x,
\quad
p_3=3.
\]

Entonces:

\[
p_1(0)=2,
\qquad
p_2(0)=1,
\qquad
p_3(0)=3.
\]

Por tanto:

\[
\boxed{
M_B(f)=
\begin{pmatrix}
4&2&6\\
2&1&3\\
6&3&9
\end{pmatrix}.
}
\]

\section{Ejercicio 9}
\subsection{Solución}

En a):

\[
A=
\begin{pmatrix}
2&-5\\
-5&4
\end{pmatrix}
\]

es simétrica.

\[
\boxed{f\text{ es simétrica}.}
\]

En b):

\[
A=
\begin{pmatrix}
0&3\\
-3&0
\end{pmatrix}
\]

es antisimétrica.

\[
\boxed{g\text{ es antisimétrica}.}
\]

En c):

\[
A=
\begin{pmatrix}
1&-2\\
7&6
\end{pmatrix}.
\]

No es simétrica ni antisimétrica.

\[
\boxed{h\text{ no es simétrica ni antisimétrica}.}
\]

\section{Ejercicio 10}
\subsection{Solución}

\[
A=
\begin{pmatrix}
2&7\\
6&-1
\end{pmatrix}.
\]

La parte simétrica es:

\[
A_s=
\frac{A+A^t}{2}
=
\begin{pmatrix}
2&\frac{13}{2}\\
\frac{13}{2}&-1
\end{pmatrix}.
\]

La parte antisimétrica es:

\[
A_a=
\frac{A-A^t}{2}
=
\begin{pmatrix}
0&\frac12\\
-\frac12&0
\end{pmatrix}.
\]

Por tanto:

\[
\boxed{
f_s(x,y)
=
2x_1y_1
+
\frac{13}{2}x_1y_2
+
\frac{13}{2}x_2y_1
-
x_2y_2
}
\]

\[
\boxed{
f_a(x,y)
=
\frac12x_1y_2-\frac12x_2y_1.
}
\]

\section{Ejercicio 11}
\subsection{Solución}

\[
P=
\begin{pmatrix}
1&1\\
-1&1
\end{pmatrix},
\qquad
Q=
\begin{pmatrix}
1&1&1\\
0&1&1\\
0&0&1
\end{pmatrix}.
\]

La nueva matriz es:

\[
M=P^tAQ.
\]

Así:

\[
\boxed{
M=
\begin{pmatrix}
-1&-6&-6\\
5&8&10
\end{pmatrix}.
}
\]

\section{Ejercicio 12}
\subsection{Solución}

La matriz de paso de la base canónica a $B$ es:

\[
P=
\begin{pmatrix}
1&3\\
2&-7
\end{pmatrix}^{-1}
=
\frac1{13}
\begin{pmatrix}
7&3\\
2&-1
\end{pmatrix}.
\]

Por congruencia:

\[
M_{B_c}(f)=P^tAP.
\]

Por tanto:

\[
\boxed{
M_{B_c}(f)
=
\frac1{13^2}
\begin{pmatrix}
150&55\\
3&18
\end{pmatrix}.
}
\]

\section{Ejercicio 13}
\subsection{Solución}

\[
A=
\begin{pmatrix}
1&1&0\\
1&2&-1\\
0&-1&1
\end{pmatrix}.
\]

\[
\boxed{\rg f=2.}
\]

Resolviendo:

\[
AX=0
\]

se obtiene:

\[
(x,y,z)=(-t,t,t).
\]

Por tanto:

\[
\boxed{
N(f)=\Span\{(-1,1,1)\}.
}
\]

La matriz:

\[
M=
\begin{pmatrix}
1&2\\
2&0
\end{pmatrix}
\]

tiene coordenadas:

\[
(1,2,0)_B.
\]

Su conjugado viene dado por:

\[
(1,2,0)AX=0.
\]

Así:

\[
\boxed{
M^c=
\{(a,b,c)_B:3a+5b-2c=0\}.
}
\]

Una base de vectores conjugados dos a dos es:

\[
\boxed{
\left\{
(1,0,0),
(0,0,1),
(-1,1,1)
\right\}_B.
}
\]

Los menores principales cumplen:

\[
\Delta_1>0,
\qquad
\Delta_2>0,
\qquad
\Delta_3=0.
\]

Por tanto:

\[
\boxed{
q\text{ es semidefinida positiva}.
}
\]

\section{Ejercicio 14}
\subsection{Solución}

La matriz asociada es:

\[
A=
\begin{pmatrix}
1&2&-3\\
2&5&-4\\
-3&-4&8
\end{pmatrix}.
\]

Su determinante es:

\[
\det A=-5\neq0.
\]

Por tanto:

\[
\rg A=3.
\]

Una matriz congruente con:

\[
\diag(1,-1,0)
\]

tendría rango $2$, pero el rango se conserva por congruencia.

Por tanto:

\[
\boxed{
\text{No existe una base con matriz }\diag(1,-1,0)
\text{ para la forma dada}.
}
\]

\section{Ejercicio 15}
\subsection{Solución}

\[
A=
\begin{pmatrix}
1&a&-1\\
a&1&2\\
-1&2&5
\end{pmatrix}.
\]

\[
\Delta_1=1,
\qquad
\Delta_2=1-a^2,
\qquad
\Delta_3=-a(5a+4).
\]

El resultado es:

\[
\boxed{
q\text{ es definida positiva si }
-\frac45<a<0
}
\]

\[
\boxed{
q\text{ es semidefinida positiva si }
a=-\frac45
\text{ o }
a=0
}
\]

y:

\[
\boxed{
q\text{ es indefinida en el resto de valores}.
}
\]

\section{Ejercicio 16}
\subsection{Solución}

\[
\Delta_1=1>0
\]

\[
\Delta_2=
\begin{vmatrix}
1&2\\
2&6
\end{vmatrix}
=2>0
\]

\[
\Delta_3
=
\det A
=
2a-2.
\]

Por tanto:

\[
\boxed{
a>1
\Longrightarrow
q\text{ definida positiva}
}
\]

\[
\boxed{
a=1
\Longrightarrow
q\text{ semidefinida positiva}
}
\]

\[
\boxed{
a<1
\Longrightarrow
q\text{ indefinida}.
}
\]

\section{Ejercicio 17}
\subsection{Solución}

La matriz es:

\[
A=
\begin{pmatrix}
1&2&1&1\\
2&3&0&1\\
1&0&-1&-1\\
1&1&-1&0
\end{pmatrix}.
\]

Una base de vectores conjugados es:

\[
v_1=(1,0,0,0)
\]

\[
v_2=(1,0,0,-1)
\]

\[
v_3=(1,0,1,-2)
\]

\[
v_4=(-1,1,0,-1).
\]

Con:

\[
P=
\begin{pmatrix}
1&1&1&-1\\
0&0&0&1\\
0&0&1&0\\
0&-1&-2&-1
\end{pmatrix},
\]

se obtiene:

\[
\boxed{
D=P^tAP
=
\diag(1,-1,2,0).
}
\]

Por tanto:

\[
\boxed{
q\text{ es indefinida y degenerada}.
}
\]

El núcleo viene dado por el vector correspondiente al elemento diagonal nulo:

\[
\boxed{
N(q)=\Span\{(-1,1,0,-1)\}.
}
\]

\section{Ejercicio 18}
\subsection{Solución}

\[
A=
\begin{pmatrix}
1&2&1\\
2&1&1\\
1&1&3
\end{pmatrix}.
\]

Una base de vectores conjugados es:

\[
v_1=(1,0,0),
\qquad
v_2=(1,-1,1),
\qquad
v_3=(4,-3,2).
\]

Entonces:

\[
P=
\begin{pmatrix}
1&1&4\\
0&-1&-3\\
0&1&2
\end{pmatrix}
\]

y:

\[
\boxed{
P^tAP
=
\diag(1,1,-7).
}
\]

Por tanto:

\[
\boxed{
q\text{ es indefinida y no degenerada}.
}
\]

\section{Ejercicio 19}
\subsection{Solución}

\[
A=
\begin{pmatrix}
1&-1&1\\
-1&0&-3\\
1&-3&-3
\end{pmatrix}.
\]

Mediante transformaciones elementales por congruencia se obtiene:

\[
\boxed{
D=
\diag(1,-1,0).
}
\]

Por tanto:

\[
\boxed{
q\text{ es indefinida y degenerada}.
}
\]

\section{Ejercicio 20}
\subsection{Solución}

Agrupando:

\[
q
=
xy+yz+xt+yt+zt
\]

\[
=
(x+z+t)(y+t)-t^2.
\]

Tomando:

\[
x+z+t=x'+y',
\qquad
y+t=x'-y',
\]

se obtiene:

\[
\boxed{
q=x'^2-y'^2-t^2.
}
\]

Aparece además una variable libre, por lo que el rango es $3$.

Por tanto:

\[
\boxed{
q\text{ es indefinida y degenerada}.
}
\]

\section{Ejercicio 21}
\subsection{Solución}

Tomamos:

\[
x=x'+z',
\qquad
z=x'-z',
\qquad
y=y'.
\]

Entonces:

\[
q
=
3y'^2+2(x'+z')(x'-z')
\]

\[
=
3y'^2+2x'^2-2z'^2.
\]

Por tanto:

\[
\boxed{
D=\diag(2,3,-2).
}
\]

Así:

\[
\boxed{
q\text{ es indefinida y no degenerada}.
}
\]

\section{Ejercicio 22}
\subsection{Solución}

\[
q
=
(x+y+z)^2
+
(a-1)(y+z)^2
+
(3-a)z^2.
\]

Por tanto:

\[
\boxed{
D=
\diag(1,a-1,3-a).
}
\]

Así:

\[
\boxed{
a<1
\Longrightarrow
q\text{ indefinida no degenerada}
}
\]

\[
\boxed{
a=1
\Longrightarrow
q\text{ semidefinida positiva}
}
\]

\[
\boxed{
1<a<3
\Longrightarrow
q\text{ definida positiva}
}
\]

\[
\boxed{
a=3
\Longrightarrow
q\text{ semidefinida positiva}
}
\]

\[
\boxed{
a>3
\Longrightarrow
q\text{ indefinida no degenerada}.
}
\]

\section{Ejercicio 23}
\subsection{Solución}

Completando cuadrados:

\[
q
=
(x+z)^2
-
a(y+2z)^2
+
(4a-1)z^2.
\]

Por tanto:

\[
\boxed{
D=
\diag(1,-a,4a-1).
}
\]

Así:

\[
\boxed{
a<0:
\text{ indefinida no degenerada}
}
\]

\[
\boxed{
a=0:
\text{ indefinida degenerada}
}
\]

\[
\boxed{
0<a<\frac14:
\text{ indefinida no degenerada}
}
\]

\[
\boxed{
a=\frac14:
\text{ indefinida degenerada}
}
\]

\[
\boxed{
a>\frac14:
\text{ indefinida no degenerada}.
}
\]

\section{Ejercicio 24}
\subsection{Solución}

La matriz es:

\[
\boxed{
A=
\begin{pmatrix}
1&0&1\\
0&1&k\\
1&k&k+1
\end{pmatrix}.
}
\]

Los menores principales son:

\[
\Delta_1=1,
\qquad
\Delta_2=1,
\qquad
\Delta_3=k(1-k).
\]

Por tanto:

\[
\boxed{
k<0\text{ o }k>1:
\text{ indefinida no degenerada}
}
\]

\[
\boxed{
0<k<1:
\text{ definida positiva}
}
\]

\[
\boxed{
k=0\text{ o }k=1:
\text{ semidefinida positiva}.
}
\]

Para $k=2$:

\[
S=\Span\{(0,0,1)\}.
\]

Entonces:

\[
(0,0,1)A
\begin{pmatrix}x\\y\\z\end{pmatrix}
=0
\]

da:

\[
\boxed{
S^c=
\{(x,y,z):x+2y+3z=0\}.
}
\]

Para $k=3$, por transformaciones de congruencia:

\[
\boxed{
D=\diag(1,1,-6)
}
\]

con:

\[
\boxed{
P=
\begin{pmatrix}
1&0&-1\\
0&1&-3\\
0&0&1
\end{pmatrix}.
}
\]

Para $k=4$, una base de vectores conjugados es:

\[
v_1=(1,0,0),
\qquad
v_2=(0,1,0),
\qquad
v_3=(-1,-4,1).
\]

Por tanto:

\[
\boxed{
P=
\begin{pmatrix}
1&0&-1\\
0&1&-4\\
0&0&1
\end{pmatrix}
}
\]

y:

\[
\boxed{
D=P^tAP=\diag(1,1,-12).
}
\]

\section{Ejercicio 25}
\subsection{Solución}

La matriz asociada es:

\[
\boxed{
A_a=
\begin{pmatrix}
1&1&2a\\
1&4&0\\
2a&0&2
\end{pmatrix}.
}
\]

Para que los vectores de $B$ sean conjugados:

\[
f(u_1,u_2)=0
\]

da:

\[
1+\lambda=0
\]

y:

\[
\lambda=-1.
\]

Además:

\[
f(u_1,u_3)=0
\]

da, como $a\neq0$:

\[
2a(1+\mu)=0
\]

y:

\[
\mu=-1.
\]

Con estos valores:

\[
f(u_2,u_3)=-5a\neq0.
\]

Por tanto:

\[
\boxed{
\text{No existen }\lambda,\mu\in\R
\text{ que hagan }B\text{ una base de vectores conjugados para }a\neq0.
}
\]

Para el carácter:

\[
\Delta_1=1
\]

\[
\Delta_2=3
\]

\[
\Delta_3
=
\det A_a
=
2(3-8a^2).
\]

Por tanto:

\[
\boxed{
q_a\text{ es definida positiva}
\iff
|a|<\sqrt{\frac38}
=
\frac{\sqrt6}{4}.
}
\]

Y:

\[
\boxed{
q_a\text{ es semidefinida positiva}
\iff
|a|=\frac{\sqrt6}{4}.
}
\]

Para $a=1$ la forma es no degenerada, por lo que para cualquier $u_3\neq0$:

\[
\dim\Span\{u_3\}^c
=
3-1.
\]

Así:

\[
\boxed{
\dim\Span\{u_3\}^c=2.
}
\]

Para $a=2$:

\[
A=
\begin{pmatrix}
1&1&4\\
1&4&0\\
4&0&2
\end{pmatrix}.
\]

\textbf{Primer método: completar cuadrados.}

\[
q
=
(x+y+4z)^2
+
3\left(y-\frac43z\right)^2
-
\frac{58}{3}z^2.
\]

Por tanto:

\[
\boxed{
D_1=
\diag\left(1,3,-\frac{58}{3}\right)
}
\]

con:

\[
\boxed{
P_1=
\begin{pmatrix}
1&-1&-\frac{16}{3}\\
0&1&\frac43\\
0&0&1
\end{pmatrix}.
}
\]

\textbf{Segundo método: base de vectores conjugados.}

Tomamos:

\[
v_1=(1,0,0),
\qquad
v_2=(-1,1,0),
\qquad
v_3=(-16,4,3).
\]

Son conjugados dos a dos y:

\[
q(v_1)=1,
\qquad
q(v_2)=3,
\qquad
q(v_3)=-174.
\]

Así:

\[
\boxed{
P_2=
\begin{pmatrix}
1&-1&-16\\
0&1&4\\
0&0&3
\end{pmatrix}
}
\]

y:

\[
\boxed{
D_2=P_2^tAP_2
=
\diag(1,3,-174).
}
\]

En ambos casos hay dos términos positivos y uno negativo.

Por tanto:

\[
\boxed{
\sg(q_2)=(2,1).
}
\]

\end{document}
